\clearpage
\begin{summarybox}
\textbf{\textcolor{CSummary}{一句话核心}}
\textbf{两切线的交点 $M$ 与弦中点 $N$，在抛物线对称轴的垂直方向上具有相同坐标；因此 $MN$ 沿着抛物线对称轴的方向。}

\medskip
\textbf{\textcolor{CSummary}{必背公式}}
对于 $y^2=2px$（$p\ne0$，$A\ne B$）：
\[
\boxed{y_M=y_N=\frac{y_1+y_2}{2}},\quad
\boxed{M_x=\frac{y_1y_2}{2p}},\quad
\boxed{N_x=\frac{y_1^2+y_2^2}{4p}}.
\]
\[
\boxed{x_N-x_M=\frac{(y_1-y_2)^2}{4p}\ne0},\quad
\boxed{|MN|=\frac{(y_1-y_2)^2}{4|p|}}.
\]
对于 $x^2=2py$：
\[
\boxed{x_M=x_N=\frac{x_1+x_2}{2}},\quad
\boxed{|MN|=\frac{(x_1-x_2)^2}{4|p|}}.
\]

\medskip
\textbf{\textcolor{CSummary}{四法优先级}}
\begin{enumerate}[leftmargin=*,itemsep=0.3em]
\item \textbf{门槛最低：}切线方程相减 $\rightarrow$ 平方差 $\rightarrow$ 中点公式。
\item \textbf{速度最快：}已知弦 $x=ty+q$ 时，直接得 $y_M=y_N=pt$。
\item \textbf{步骤最少：}切点弦 $y_My=p(x+x_M)$ $\rightarrow$ 两根之和。
\item \textbf{常规通用：}完整求出 $M,N$ 的坐标，再验证水平方向的距离非零。
\end{enumerate}

\medskip
\textbf{\textcolor{CSummary}{几何记忆线索}}
\textbf{纵坐标取平均 $\rightarrow$ 同一高度 $\rightarrow$ 水平连线。}
若共同纵坐标不为 $0$，与 $x$ 轴平行；若为 $0$，就与 $x$ 轴重合。
另一个标准形式只需交换横、纵坐标。

\medskip
\textbf{\textcolor{CSummary}{考前检查}}
\begin{itemize}[leftmargin=*,itemsep=0.2em,topsep=0.3em]
\item $A,B$ 必须是不同点；$p\ne0$；并区分 $p$ 与 $2p$。
\item 顶点切线可用 $yy_0=p(x+x_0)$，不要强行除以 $y_0$。
\item 计算长度必须用 $|p|$；对称切点可能使 $MN$ 与坐标轴重合。
\end{itemize}
\end{summarybox}
